\chapter{迴圈}
\label{ch:loop}

到目前為止，我們寫的程式多半只執行一次就結束。 讀一個數字、算一個結果、印出
來。 但實際上很多問題需要\textbf{重複}做同一件事：印出 $N$ 行相同的文字、把
一串數字逐一讀進來、一直問使用者輸入直到格式正確、或是模擬每一天細菌怎麼繁殖。
如果每一輪都手動複製貼上同一段程式碼，程式不但冗長，改一個地方還得改很多
處，幾乎無法維護。

\textbf{迴圈}（loop）就是為了解決這個問題：把「要重複做的事」寫成一段，再告
訴電腦「重複幾次」或「重複到什麼條件為止」。 C 語言提供三種迴圈語法。 
\texttt{for}、\texttt{while}、\texttt{do-while}。 語法不同，但背後都是同一
套機制：每次迭代（iteration）更新迴圈變數或狀態，直到條件不再成立為止。 從
硬體的角度看，CPU 其實只是在反覆執行「比較條件 → 跳回或繼續往下」這兩步；
迴圈讓我們用高階語法描述這個過程，而不必自己寫跳轉指令。

這一章依語法複雜度推進：先認識三種迴圈與 \texttt{for}/\texttt{while} 的等價
關係，再學 \texttt{break} 與 \texttt{continue}，接著用巢狀迴圈輸出圖形，然
後練解題平台必備的輸入模板，最後才進入數論算法與需要追蹤多個狀態的模擬題。
題庫裡這個單元有八十多題，難度跨度很大；我們沿著整條難度曲線走，把核心語法
與經典解題模型分開講清楚。

\section{固定次數的重複：for 迴圈}
\label{sec:loop-for-basics}

\texttt{for} 迴圈最適合「次數事先知道」的情形：從 $0$ 數到 $n-1$、從 $1$ 數
到 $n$、或是倒數 $n$ 次。 它的結構把「初始化」「條件」「每次迭代後的更新」
集中寫在一行，一眼就能看出迴圈跑幾輪。

\begin{code}{c}{for 迴圈的基本形狀}
for (初始化; 條件; 更新) {
    重複執行的敘述;
}
\end{code}

範例 \thechapter.2 讀入正整數 $n$，把同一行文字印 $n$ 次。 這是最單純的
「已知次數、每次做同一件事」的模型。

\begin{example}{c}{(print-fixed-line.c) 把固定文字印 N 次}
#include <stdio.h>

int main(void) {
    int n;
    scanf("%d", &n);

    for (int i = 0; i < n; i++) {
        printf("Hello\n");
    }
    return 0;
}
\end{example}

\iosample{
3 \\
}{
Hello \\
Hello \\
Hello \\
}

\texttt{for (int i = 0; i < n; i++)} 裡，\texttt{i} 從 $0$ 開始，每輪加
$1$，當 \texttt{i == n} 時條件不成立、迴圈結束，所以總共執行 $n$ 次。
這跟陣列索引從 $0$ 開始的慣例一致。 之後用 \texttt{for} 走訪陣列時，幾
乎都會看到 \texttt{i < n} 這種寫法。


\keypoint{\texttt{for (int i = 0; ...)} 裡宣告的 \texttt{i}，作用域就是這
一句 \texttt{for}；迴圈結束後在 \texttt{main} 裡不能再寫 \texttt{printf("\%d", i)}。
若要在迴圈外使用計數器，把 \texttt{int i} 提到 \texttt{for} 前面。}

\keypoint{\texttt{for} 迴圈的三個部分（初始化、條件、更新）都在迴圈\textbf{開
始前}就決定好；條件在每一輪\textbf{一開始}先檢查，不成立就整個跳過，一次都
不執行。}

\begin{drills}[label=drills:loop-for-basics]
\item \texttt{for (int i = 0; i < n; i++)} 印 \texttt{Hello} $n$ 次，
若輸入 $n = 0$，迴圈本體會執行幾次？
\item \texttt{for} 裡宣告的 \texttt{int i}，離開 \texttt{for} 這一句之後
還能不能在 \texttt{main} 裡繼續用？
\item \texttt{i < n} 與 \texttt{i <= n} 兩種條件，各自讓迴圈跑幾次？
\end{drills}

\section{改變計數方向：倒數與反向列印}
\label{sec:loop-countdown}

次數固定，但「每一輪要做的事」跟計數器的值有關時，只要改初始化、條件、更新
的方向就好。 範例 \thechapter.3 從 $n$ 倒數印到 $1$，每一行印一個數字。

\begin{example}{c}{(print-backward.c) 從 n 倒數印到 1}
#include <stdio.h>

int main(void) {
    int n;
    scanf("%d", &n);

    for (int i = n; i >= 1; i--) {
        printf("%d\n", i);
    }
    return 0;
}
\end{example}

\iosample{
5 \\
}{
5 \\
4 \\
3 \\
2 \\
1 \\
}

範例 \thechapter.4 則把 $n$ 到 $0$ 全部印在同一行，數字之間以空白分隔。 
這裡多了一個常見的輸出細節：最後一個數字後面不要多印空白。

\begin{example}{c}{(countdown.c) 倒數 n 到 0，空白分隔}
#include <stdio.h>

int main(void) {
    int n;
    scanf("%d", &n);

    for (int i = n; i >= 0; i--) {
        printf("%d", i);
        if (i > 0) {
            printf(" ");
        }
    }
    printf("\n");
    return 0;
}
\end{example}

\iosample{
3 \\
}{
3 2 1 0 \\
}

以 $n = 3$ 走一遍：\texttt{i} 依序是 $3, 2, 1, 0$；印完 $3$、$2$、$1$ 時
\texttt{i > 0} 成立，後面補空白；印到 $0$ 時不再補空白，最後換行。

\stylepoint{輸出一排以空白分隔的值時，用 \texttt{if (i > 0)}（或類似判
斷）避免行尾多餘空白，比事後再處理字串乾淨得多。}

\begin{drills}[label=drills:loop-countdown]
\item \texttt{for (int i = n; i >= 1; i--)} 跟
\texttt{for (int i = n; i >= 0; i--)}，哪一個會多印一個數字？多印的是哪
一個？
\item \texttt{countdown.c} 為什麼要用 \texttt{if (i > 0)} 才印空白，而
不是每個數字後面都印？
\item 若把 \texttt{if (i > 0)} 誤寫成 \texttt{if (i >= 0)}，輸出會有什
麼問題？
\end{drills}

\section{while 迴圈：條件成立就繼續}
\label{sec:loop-while-basics}

\texttt{while} 迴圈適合「不知道要跑幾次，只知道什麼時候該停」的情形。 它只
有條件，沒有把初始化跟更新寫在括號裡。 這些通常要自己在迴圈本體裡或迴圈
前後處理。

\begin{code}{c}{while 迴圈的基本形狀}
while (條件) {
    重複執行的敘述;
}
\end{code}

範例 \thechapter.6 讀入正整數 $n$，只要 $n$ 還是偶數就除以 $2$，直到變
成奇數為止。 這種「一直做某個操作直到狀態改變」的模式，用 \texttt{while}
比 \texttt{for} 自然。

\begin{example}{c}{(divide-until-odd.c) 一直除以 2 直到變成奇數}
#include <stdio.h>

int main(void) {
    int n;
    scanf("%d", &n);

    while (n % 2 == 0) {
        n /= 2;
    }
    printf("%d\n", n);
    return 0;
}
\end{example}

\iosample{
24 \\
}{
3 \\
}

$n = 24$ 時：$24 \to 12 \to 6 \to 3$，三次除以 $2$ 後 $3$ 是奇數，迴圈
結束，印出 $3$。


若本體裡忘了更新條件相關的變數，迴圈可能永遠停不下來：

\begin{code}{c}{危險：忘記更新，形成無窮迴圈}
int i = 0;
while (i < n) {
    printf("%d\n", i);
    // 漏寫 i++：i 永遠是 0，條件永遠成立
}
\end{code}

終端機若卡住不動、風扇開始狂轉，先檢查迴圈裡有沒有把計數器或狀態往終止
條件推進。 開發時用 \texttt{Ctrl-C} 強制結束程式，再回去補上漏掉的更新。

\keypoint{\texttt{while} 迴圈在每一輪\textbf{一開始}先檢查條件；如果一開始
條件就不成立，迴圈本體一次都不會執行。 \texttt{do-while} 則把條件
檢查放在本體之後，至少會跑一次；後面會專門對比。}

\begin{drills}[label=drills:loop-while-basics]
\item \texttt{divide-until-odd.c} 輸入 $24$，迴圈總共跑幾輪？每一輪
\texttt{n} 依序變成什麼？
\item 範例裡若漏寫 \texttt{n /= 2;}，會發生什麼事？
\item \texttt{while} 迴圈在每一輪的\emph{開始}還是\emph{結束}檢查條
件？
\end{drills}

\section{for 跟 while 其實是同一件事}
\label{sec:loop-for-while-equiv}

\texttt{for (初始化; 條件; 更新)} 可以拆成等價的 \texttt{while}：初始
化只做一次，然後「先測條件、再執行本體、最後更新」，再回頭測條件。
條件一開始就不成立時，本體一次都不會跑。 這點 \texttt{for} 與
\texttt{while} 相同，跟 \texttt{do-while} 不同。

\begin{code}{c}{把 for 拆成等價的 while}
初始化;
while (條件) {
    本體;
    更新;
}
\end{code}

範例 \thechapter.9 用 \texttt{while} 印出 $0$ 到 $n-1$，效果等同
\texttt{for (int i = 0; i < n; i++)}。 輸入 \texttt{3} 得到三行
\texttt{0}、\texttt{1}、\texttt{2}。

\begin{example}{c}{(for-as-while.c) 用 while 模擬固定次數的 for}
#include <stdio.h>

int main(void) {
    int n;
    scanf("%d", &n);
    int i = 0;
    while (i < n) {
        printf("%d\n", i);
        i++;
    }
    return 0;
}
\end{example}

\iosample{
3 \\
}{
0 \\
1 \\
2 \\
}

\begin{center}
\begin{tikzpicture}[node distance=8mm and 20mm]
    \node[process] (init) {初始化};
    \node[decision, below=of init] (cond) {條件？};
    \node[process, below=of cond] (body) {本體};
    \node[process, right=of body] (upd) {更新};
    \node[process, below=of body] (done) {離開迴圈};
    \draw[flowarrow] (init.south) -- (cond.north);
    \draw[flowarrow] (cond.south) -- node[left,font=\scriptsize] {是} (body.north);
    \draw[flowarrow] (body.east) -- (upd.west);
    \draw[flowarrow] (upd.north) |- (cond.east);
    \draw[flowarrow] (cond.west) -- ++(-2cm,0) |- node[below,font=\scriptsize,pos=0.25] {否} (done.west);
\end{tikzpicture}
\diagcaption{\texttt{for} 與 \texttt{while} 的共同流程：先測條件。 條
件為假時本體一次也不執行。}
\label{fig:ch05-for-while-flow}
\end{center}

計數器作用域已在第一節說明；這裡專心看 \texttt{for} 怎麼改寫成
\texttt{while}，以及兩者共用的流程。

差一錯誤（off-by-one）幾乎都出在條件用 \texttt{<} 還是 \texttt{<=}。
\texttt{i = 0; i < n} 跑 $n$ 次；\texttt{i = 1; i <= n} 也跑 $n$
次。 兩種都對，混用就容易多跑或少跑一輪。 先在紙上寫出 $i$ 會取哪些
值，再寫條件。

逗號運算子可以把兩個運算式合成一個，值是右邊那個。 偶爾會看到
\texttt{for (i = 0, sum = 0; i < n; i++)} 這種「初始化兩個變數」的寫
法；能讀懂即可，不必為了短而把更新也擠進逗號裡。

\begin{code}{c}{for 的初始化裡用逗號隔開兩個指定}
for (i = 0, sum = 0; i < n; i++) {
    sum += i;
}
\end{code}

\begin{drills}[label=drills:loop-for-while-equiv]
\item 把 \texttt{for (初始化; 條件; 更新)} 拆成等價的 \texttt{while}
時，「更新」要放在迴圈本體的最前面還是最後面？
\item \texttt{for} 與 \texttt{while}，當條件一開始就不成立時，本體各
會執行幾次？
\item \texttt{for-as-while.c} 輸入 $3$，會印出哪三個數字？
\end{drills}

\section{do-while：至少執行一次}
\label{sec:loop-do-while}

\texttt{do-while} 把條件檢查放在迴圈\textbf{本體之後}，所以不管條件一開始
成不成立，本體至少會跑一次。 這在「一定要先問一次、再決定要不要繼續問」
的輸入驗證特別有用。

\begin{code}{c}{do-while 迴圈的基本形狀}
do {
    重複執行的敘述;
} while (條件);
\end{code}

範例 \thechapter.12 一直讀整數，直到讀到嚴格正數為止。 使用者可能第一次就
輸入錯的，但我們仍然要先讀進來才能判斷。

\begin{example}{c}{(read-until-positive.c) 讀到嚴格正數為止}
#include <stdio.h>

int main(void) {
    int x;

    do {
        scanf("%d", &x);
    } while (x <= 0);

    printf("%d\n", x);
    return 0;
}
\end{example}

\iosample{
-1 \\
0 \\
5 \\
}{
5 \\
}

第一次讀到 $-1$、第二次 $0$，條件 \texttt{x <= 0} 都成立，繼續讀；第三
次讀到 $5$，條件不成立，跳出迴圈。

\keypoint{\texttt{do-while} 保證迴圈本體至少執行一次；輸入驗證、選單、
「重試直到成功」這類題目，幾乎都適合用 \texttt{do-while} 而不是
\texttt{while}。}

\texttt{while} 與 \texttt{for} 都可能一次都不進本體。 \texttt{do-while}
相反：先做，再問。 範例 \thechapter.13 把 \texttt{n} 設成 $5$，迴圈條
件卻是 \texttt{n < 0}。 條件從頭到尾都是假，但仍會印出一行 \texttt{5}，
因為本體在測試之前就跑過了。

\begin{example}{c}{(do-once.c) 條件為假時，do-while 仍執行一次}
#include <stdio.h>

int main(void) {
    int n = 5;
    do {
        printf("%d\n", n);
        n++;
    } while (n < 0);
    return 0;
}
\end{example}

\iosample{}{
5 \\
}

這就是為什麼輸入驗證愛用 \texttt{do-while}：你總得先讀一次，才知道合
不合法。

\begin{drills}[label=drills:loop-do-while]
\item \texttt{do-while} 跟 \texttt{while} 最大的差別是什麼？
\item \texttt{do-once.c} 裡 \texttt{n = 5}，迴圈條件是
\texttt{n < 0}（從頭到尾都是假），但程式仍會印出什麼？為什麼？
\item \texttt{read-until-positive.c} 若第一個輸入就是正數，迴圈本體
會執行幾次？
\end{drills}

\section{break、continue 與提早離開迴圈}
\label{sec:loop-break-continue}

有時候迴圈不必跑完所有輪次。 找到答案就可以 \texttt{break} 跳出；某些
輪次要跳過則用 \texttt{continue} 直接進下一輪。 後面的質數判定與
Collatz 序列都會再用到 \texttt{break}。

\begin{code}{c}{在 for 迴圈裡用 break 提早結束}
for (int i = 0; i < n; i++) {
    if (找到答案) {
        break;
    }
}
\end{code}


範例 \thechapter.15 讀入 $n$，從 $1$ 掃到 $n$，找到第一個 $7$ 的倍數就
\texttt{break}，印出那個數；若都沒有就印 \texttt{-1}。

\begin{example}{c}{(first-multiple-of-7.c) 找到第一個 7 的倍數就停}
#include <stdio.h>

int main(void) {
    int n;
    scanf("%d", &n);

    int found = -1;
    for (int i = 1; i <= n; i++) {
        if (i % 7 == 0) {
            found = i;
            break;
        }
    }
    printf("%d\n", found);
    return 0;
}
\end{example}

\iosample{
10 \\
}{
7 \\
}

\iosample{
5 \\
}{
-1 \\
}

\begin{code}{c}{在 for 迴圈裡用 continue 跳過本輪}
for (int i = 0; i < n; i++) {
    if (這一輪要跳過) {
        continue;
    }
    處理本輪;
}
\end{code}

\stylepoint{一個迴圈盡量只有一個入口、一個出口，可讀性較高；但「找到就
停」這類明確的 \texttt{break}，通常比設一堆旗標變數更直覺。 不要為了
避免 \texttt{break} 而把邏輯寫繞。}

\texttt{continue} 適合「這一輪剩下的事可以整段不理」。 範例
\thechapter.17 讀 $n$ 個整數，只把正數加進總和；非正數一碰到就
\texttt{continue}，後面的 \texttt{sum += x} 不會執行。 輸入
\texttt{4} 筆資料 \texttt{7 -2 3 8}，略過 $-2$，得到 $18$。

\begin{example}{c}{(sum-positives.c) 用 continue 跳過非正數}
#include <stdio.h>

int main(void) {
    int n;
    scanf("%d", &n);
    int sum = 0;
    for (int count = 0; count < n; count++) {
        int x;
        scanf("%d", &x);
        if (x <= 0) {
            continue;
        }
        sum += x;
    }
    printf("%d\n", sum);
    return 0;
}
\end{example}

\iosample{
4 7 -2 3 8 \\
}{
18 \\
}

同一題也可以寫成「只有正數才加」，如範例 \thechapter.18。 結果一樣，
但讀程式的人看到 \texttt{if (x > 0)} 往往會往下找 \texttt{else}。 其
實沒有 else，只是「其餘情況什麼都不做」。 \texttt{continue} 把「不理
這一輪」寫在條件成立的當下，後面的加總邏輯就專心處理正數。

\begin{example}{c}{(sum-positives-if.c) 不用 continue 的等價寫法}
#include <stdio.h>

int main(void) {
    int n;
    scanf("%d", &n);
    int sum = 0;
    for (int count = 0; count < n; count++) {
        int x;
        scanf("%d", &x);
        if (x > 0) {
            sum += x;
        }
    }
    printf("%d\n", sum);
    return 0;
}
\end{example}

\iosample{
4 7 -2 3 8 \\
}{
18 \\
}

迴圈本體裡宣告的 \texttt{x}，離開那一對大括號就不能再用。 這正好符合
「它只是這一輪的暫存」。 宣告和使用靠在一起，比把 \texttt{x} 提到
\texttt{main} 開頭更清楚。

\keypoint{\texttt{continue} 結束的是\textbf{這一輪}，不是整個迴圈；
\texttt{break} 才是整段離開。 兩者都不要拿來代替本來該寫清楚的條件。}

\begin{drills}[label=drills:loop-break-continue]
\item \texttt{break} 跟 \texttt{continue}，哪一個是整段離開迴圈，哪
一個只結束這一輪？
\item \texttt{first-multiple-of-7.c} 輸入 $5$（找 $1$ 到 $5$ 裡第一個
$7$ 的倍數），會印出什麼？
\item \texttt{sum-positives.c} 輸入 \texttt{4 7 -2 3 8}，\texttt{continue}
跳過的是哪一個數字？最後 \texttt{sum} 是多少？
\end{drills}

\section{兩段式輸出：先上後下}
\label{sec:loop-two-stage}

有些題目要依序做兩段不同但結構類似的輸出。 與其寫一個複雜的迴圈，不如老老
實實寫兩個 \texttt{for}：第一段 $1$ 到 $n$，第二段 $n-1$ 到 $1$。 可讀性
通常比硬塞在一個迴圈裡好。

\begin{example}{c}{(up-and-down.c) 先印 1 到 n，再印 n-1 到 1}
#include <stdio.h>

int main(void) {
    int n;
    scanf("%d", &n);

    for (int i = 1; i <= n; i++) {
        printf("%d\n", i);
    }
    for (int i = n - 1; i >= 1; i--) {
        printf("%d\n", i);
    }
    return 0;
}
\end{example}

\iosample{
3 \\
}{
1 \\
2 \\
3 \\
2 \\
1 \\
}

兩個迴圈各自獨立，第二段從 \texttt{n - 1} 開始而不是 \texttt{n}，避免
$n$ 被印兩次。 這種「邊界少印或多印一次」是迴圈題目最常見的失誤之一。

\stylepoint{兩段結構相似的輸出，用兩個清楚的迴圈分開寫，通常比一個迴圈夾
雜 \texttt{if} 判斷「現在是上升還是下降段」更容易讀、也比較不容易算錯邊
界。}

\begin{drills}[label=drills:loop-two-stage]
\item 範例 \texttt{up-and-down.c} 裡，第二個 \texttt{for} 為什麼要從
\texttt{n - 1} 開始，而不是從 \texttt{n} 開始？
\item 如果兩段都從 \texttt{n} 開始遞減／遞增，$n = 3$ 的輸出會多印出
什麼？
\end{drills}

\section{巢狀迴圈：矩形與數字圖形}
\label{sec:loop-nested-shapes}

當外層迴圈控制「第幾列」、內層控制「這一列印幾次或印什麼」，就形成\textbf{
巢狀迴圈}（nested loop）。 內層每跑完一輪，外層才前進一格。 總執行次數
大約是外層次數乘以內層次數。

\begin{code}{c}{巢狀 for 迴圈印矩形}
for (int row = 0; row < height; row++) {
    for (int col = 0; col < width; col++) {
        printf("#");
    }
    printf("\n");
}
\end{code}

範例 \thechapter.21 讀入寬 $W$、高 $H$，印出 $H$ 行、每行 $W$ 個
\texttt{\#}。

\begin{example}{c}{(print-rectangle.c) 印 W x H 的井號矩形}
#include <stdio.h>

int main(void) {
    int w, h;
    scanf("%d %d", &w, &h);

    for (int row = 0; row < h; row++) {
        for (int col = 0; col < w; col++) {
            printf("#");
        }
        printf("\n");
    }
    return 0;
}
\end{example}

\iosample{
4 3 \\
}{
\#\#\#\# \\
\#\#\#\# \\
\#\#\#\# \\
}

範例 \thechapter.22 第 $i$ 列印數字 $i$ 共 $i$ 次。 內層次數跟外層變數
綁在一起，形成三角形。

\begin{example}{c}{(number-triangle.c) 第 i 列重複數字 i 共 i 次}
#include <stdio.h>

int main(void) {
    int n;
    scanf("%d", &n);

    for (int i = 1; i <= n; i++) {
        for (int j = 0; j < i; j++) {
            printf("%d", i);
        }
        printf("\n");
    }
    return 0;
}
\end{example}

\iosample{
4 \\
}{
1 \\
22 \\
333 \\
4444 \\
}

範例 \thechapter.23 第 $i$ 列印 $1$ 到 $i$。 同一個外層 \texttt{i}，
內層從 $1$ 數到 \texttt{i}，圖形從「重複同一數字」變成「遞增序列」。

\begin{example}{c}{(repeated-pattern.c) 第 i 列印 1 到 i}
#include <stdio.h>

int main(void) {
    int n;
    scanf("%d", &n);

    for (int i = 1; i <= n; i++) {
        for (int j = 1; j <= i; j++) {
            printf("%d", j);
        }
        printf("\n");
    }
    return 0;
}
\end{example}

\iosample{
4 \\
}{
1 \\
12 \\
123 \\
1234 \\
}

\stylepoint{巢狀迴圈寫圖形時，先想清楚「外層控制列還是行、內層控制這一列要
印幾個字元或什麼內容」，在紙上畫小表格對一下，比直接寫程式瞎試省時。}

\begin{drills}[label=drills:loop-nested-shapes]
\item 巢狀迴圈裡，外層通常控制什麼、內層通常控制什麼？
\item \texttt{number-triangle.c} 第 $i$ 列印數字 $i$ 共幾次？$n = 4$
時第 $3$ 列印出什麼？
\item \texttt{repeated-pattern.c} 跟 \texttt{number-triangle.c} 印出的
圖形有什麼不同？
\end{drills}

\section{乘法表與週期圖樣}
\label{sec:loop-mult-table}

$N \times N$ 乘法表是巢狀迴圈的經典練習：外層 \texttt{i}、內層 \texttt{j}，
印 \texttt{i * j}，並用 \texttt{\%2d} 對齊欄寬。

\begin{example}{c}{(multiplication-table.c) N x N 乘法表}
#include <stdio.h>

int main(void) {
    int n;
    scanf("%d", &n);

    for (int i = 1; i <= n; i++) {
        for (int j = 1; j <= n; j++) {
            printf("%2d", i * j);
        }
        printf("\n");
    }
    return 0;
}
\end{example}

\iosample{
3 \\
}{
 1 2 3 \\
 2 4 6 \\
 3 6 9 \\
}

範例 \thechapter.25 在 $N \times N$ 網格上，行列索引之和能被 $3$ 整除的
格子印 \texttt{@}，其餘印 \texttt{-}。 這類「週期圖樣」題目，關鍵往往是
找到 \texttt{(i + j) \% k} 或類似規律。

\begin{example}{c}{(diagonal-tile.c) 週期對角線圖樣}
#include <stdio.h>

int main(void) {
    int n;
    scanf("%d", &n);

    for (int i = 0; i < n; i++) {
        for (int j = 0; j < n; j++) {
            if ((i + j) % 3 == 0) {
                printf("@");
            } else {
                printf("-");
            }
        }
        printf("\n");
    }
    return 0;
}
\end{example}

\iosample{
4 \\
}{
@--@ \\
--@- \\
-@-- \\
@--@ \\
}

\begin{drills}[label=drills:loop-mult-table]
\item \texttt{multiplication-table.c} 裡 \texttt{printf("\%2d", ...)}
的 \texttt{\%2d} 是做什麼用的？
\item \texttt{diagonal-tile.c} 用 \texttt{(i + j) \% 3 == 0} 判斷要印
\texttt{@} 還是 \texttt{-}；$n = 4$ 時第一列（\texttt{i = 0}）有哪幾欄
是 \texttt{@}？
\end{drills}

\section{置中的數字金字塔}
\label{sec:loop-pyramid}

巢狀迴圈不一定印實心矩形。 有時內層要先印「左邊的空白（或底線）」，
再印這一列的內容。 範例 \thechapter.26 第 $i$ 列先印 $n-i$ 個底線，再
把數字 $i$ 重複 $i$ 次。 $n = 3$ 時外形會像金字塔。

\begin{example}{c}{(number-pyramid.c) 用底線對齊的數字金字塔}
#include <stdio.h>

int main(void) {
    int n;
    scanf("%d", &n);
    for (int i = 1; i <= n; i++) {
        for (int s = 0; s < n - i; s++) {
            printf("_");
        }
        for (int k = 0; k < i; k++) {
            printf("%d", i);
        }
        printf("\n");
    }
    return 0;
}
\end{example}

\iosample{
3 \\
}{
\_\_1 \\
\_22 \\
333 \\
}

外層仍管「第幾列」，內層拆成兩段：一段專管左邊界、一段專管內容。 先
在紙上把 $n=3$ 的三列畫出來，再數底線個數，比直接在終端機裡試縮排快。

\begin{drills}[label=drills:loop-pyramid]
\item \texttt{number-pyramid.c} 第 $i$ 列先印幾個底線、再印幾個數
字？
\item 內層的兩段 \texttt{for} 分別負責圖形的哪一部分？
\end{drills}

\section{do-while 驗證輸入：密碼、分數與字元}
\label{sec:loop-input-validation}

同一套「讀入 → 檢查 → 不對就再讀」可以套在不同型別的輸入上。 範例
\thechapter.27 模擬 PIN 碼驗證：不是 \texttt{1234} 就一直重試。

\begin{example}{c}{(pin-until-match.c) 重試直到 PIN 正確}
#include <stdio.h>

int main(void) {
    int pin;

    do {
        scanf("%d", &pin);
    } while (pin != 1234);

    printf("OK\n");
    return 0;
}
\end{example}

\iosample{
9999 \\
1234 \\
}{
OK \\
}

範例 \thechapter.28 限制分數在 $0$ 到 $100$ 之間。 條件變成區間檢查，
結構完全一樣。

\begin{example}{c}{(valid-grade.c) 讀到 0 到 100 的分數為止}
#include <stdio.h>

int main(void) {
    int score;

    do {
        scanf("%d", &score);
    } while (score < 0 || score > 100);

    printf("%d\n", score);
    return 0;
}
\end{example}

\iosample{
120 \\
-5 \\
85 \\
}{
85 \\
}

讀字元時，空白字元（空格、換行）常會留在輸入緩衝區裡干擾下一次
\texttt{scanf}。 在 \texttt{\%c} 前面加一個空格。 寫成 \texttt{" \%c"}。 
會讓 \texttt{scanf} 先跳過前面的空白再讀一個字元。 範例 \thechapter.29
用這個技巧讀 \texttt{Y} 或 \texttt{N}。

\begin{example}{c}{(yes-or-no.c) 讀到 Y 或 N 為止}
#include <stdio.h>

int main(void) {
    char c;

    do {
        scanf(" %c", &c);
    } while (c != 'Y' && c != 'N');

    printf("%c\n", c);
    return 0;
}
\end{example}

\iosample{
A \\
M \\
Y \\
}{
Y \\
}

\stylepoint{用 \texttt{scanf(" \%c", \&c)} 讀單一字元時，格式字串開頭的空
格會跳過換行與空白，避免上一次 \texttt{scanf} 留下的換行被當成輸入讀進
來。}

範例 \thechapter.30 讀星期（$1$ 到 $7$），不在範圍就再讀，直
到合法為止。 輸入 \texttt{0 8 3} 時前兩次被丟掉，最後印 \texttt{3}。

\begin{example}{c}{(valid-weekday.c) 讀到 1 到 7 的星期為止}
#include <stdio.h>

int main(void) {
    int d;
    do {
        scanf("%d", &d);
    } while (d < 1 || d > 7);
    printf("%d\n", d);
    return 0;
}
\end{example}

\iosample{
0 8 3 \\
}{
3 \\
}
\begin{drills}[label=drills:loop-input-validation]
\item 這一節的每個例子都用 \texttt{do-while} 而不是 \texttt{while}
驗證輸入，為什麼？
\item \texttt{scanf(" \%c", \&c)} 格式字串開頭的空格，作用是什麼？
\item \texttt{valid-grade.c} 輸入依序是 \texttt{120}、\texttt{-5}、
\texttt{85}，迴圈本體執行幾次？最後印出什麼？
\end{drills}

\section{哨兵值、選單與累加}
\label{sec:loop-sentinel}

\textbf{哨兵值}（sentinel）是一個約定好的特殊輸入，代表「資料結束了」。 
例如 $-1$ 或 $0$。 搭配 \texttt{do-while} 或 \texttt{while}，可以讀一
串未知長度的資料並做統計。

範例 \thechapter.31 模擬選單：選 $0$ 離開，其他選項印 \texttt{AGAIN}。

\begin{example}{c}{(menu-until-quit.c) 選單迴圈直到選 0}
#include <stdio.h>

int main(void) {
    int choice;

    do {
        scanf("%d", &choice);
        if (choice != 0) {
            printf("AGAIN\n");
        }
    } while (choice != 0);

    printf("BYE\n");
    return 0;
}
\end{example}

\iosample{
1 \\
2 \\
0 \\
}{
AGAIN \\
AGAIN \\
BYE \\
}

範例 \thechapter.32 在讀到 $-1$ 之前，數有幾筆有效資料（$-1$ 本身不計）。

\begin{example}{c}{(sum-until-sentinel.c) 以 -1 為哨兵值計數}
#include <stdio.h>

int main(void) {
    int x, count = 0;

    do {
        scanf("%d", &x);
        if (x != -1) {
            count++;
        }
    } while (x != -1);

    printf("%d\n", count);
    return 0;
}
\end{example}

\iosample{
1 \\
2 \\
3 \\
-1 \\
}{
3 \\
}

範例 \thechapter.33 則把哨兵改成 $0$，讀到的所有數字（包含最後那個 $0$）
加總。 注意這裡 $0$ 既參與累加，又觸發結束。

\begin{example}{c}{(sum-until-zero.c) 累加直到讀到 0}
#include <stdio.h>

int main(void) {
    int x, sum = 0;

    do {
        scanf("%d", &x);
        sum += x;
    } while (x != 0);

    printf("%d\n", sum);
    return 0;
}
\end{example}

\iosample{
1 \\
2 \\
3 \\
0 \\
}{
6 \\
}

\keypoint{哨兵值要不要參與計算，完全看題目規定；寫程式前先確認「結束符號
本身算不算一筆資料」，比寫完再改邏輯省時。}

\begin{drills}[label=drills:loop-sentinel]
\item 什麼是「哨兵值」？
\item \texttt{sum-until-sentinel.c}（以 $-1$ 為哨兵）跟
\texttt{sum-until-zero.c}（以 $0$ 為哨兵）最大的差別是什麼？
\item \texttt{menu-until-quit.c} 輸入 \texttt{1 2 0}，會依序印出哪幾
行？
\end{drills}

\section{多筆測資：已知筆數與 Case 前綴}
\label{sec:loop-multi-case}

競賽與題庫常見的輸入格式是：第一行給測資筆數 $N$，後面接 $N$ 組資料。
這時外層用 \texttt{for (int i = 0; i < n; i++)} 包住內層處理即可。

\begin{example}{c}{(multi-case-compare.c) N 組資料比較兩數大小}
#include <stdio.h>

int main(void) {
    int n;
    scanf("%d", &n);

    for (int i = 0; i < n; i++) {
        int a, b;
        scanf("%d %d", &a, &b);
        if (a > b) {
            printf("A\n");
        } else if (a < b) {
            printf("B\n");
        } else {
            printf("T\n");
        }
    }
    return 0;
}
\end{example}

\iosample{
3 \\
1 2 \\
2 2 \\
3 1 \\
}{
B \\
T \\
A \\
}

有些題目要求每組輸出前面加 \texttt{Case k:} 前綴。 範例 \thechapter.35
把迴圈計數器從 $1$ 開始（\texttt{k = 1; k <= n; k++}），直接當 Case 編
號。

\begin{example}{c}{(case-prefix.c) 多筆測資加上 Case 前綴}
#include <stdio.h>

int main(void) {
    int n;
    scanf("%d", &n);

    for (int k = 1; k <= n; k++) {
        int a, b;
        scanf("%d %d", &a, &b);
        printf("Case %d: ", k);
        if (a > b) {
            printf("A\n");
        } else if (a < b) {
            printf("B\n");
        } else {
            printf("T\n");
        }
    }
    return 0;
}
\end{example}

\iosample{
2 \\
1 2 \\
2 2 \\
}{
Case 1: B \\
Case 2: T \\
}

測資筆數未知、直到 EOF（檔案結束）為止時，可以用 \texttt{while (scanf(...)
!= EOF)} 這種寫法。 檔案處理那一章會更完整地讲，這裡先記住「外層迴圈控
制還有沒有下一組資料」這個概念就好。

\begin{code}{c}{一直讀到 EOF 為止}
while (scanf("%d %d", &a, &b) != EOF) {
    處理這一組資料;
}
\end{code}

另一種常見終止條件是「某個特殊值組合」。 例如兩個 \texttt{0} 代表結束。 
邏輯跟哨兵值相同，只是形式從單一數字變成多個欄位。


範例 \thechapter.37 是解題平台上最常見的「測資筆數未知、讀到 EOF 為止」模板：
每讀到兩個整數就印它們的和，沒有下一行時 \texttt{scanf} 回傳 \texttt{EOF}、
迴圈結束。 檔案與標準輸入重定向的細節留給後面章節；這裡先記住這個外層
\texttt{while} 骨架。

\begin{example}{c}{(sum-until-eof.c) 一直讀兩整數相加，直到 EOF}
#include <stdio.h>

int main(void) {
    int a, b;
    while (scanf("%d %d", &a, &b) != EOF) {
        printf("%d\n", a + b);
    }
    return 0;
}
\end{example}

\iosample{
1 2 \\
3 4 \\
}{
3 \\
7 \\
}


\begin{drills}[label=drills:loop-multi-case]
\item 「先給測資筆數 $N$，再接 $N$ 組資料」這種輸入格式，通常用什麼
迴圈結構處理？
\item \texttt{case-prefix.c} 跟 \texttt{multi-case-compare.c} 比起
來，多做了什麼事？
\item 測資筆數\emph{未知}、要讀到檔案結束為止時，可以用哪種
\texttt{while} 寫法？
\end{drills}

\section{連續整數求和}
\label{sec:loop-range-sum}

求 $A + (A+1) + \cdots + B$ 這類問題，用迴圈從 $A$ 加到 $B$ 最直接。
如果輸入可能 $A > B$，先交換再算。 這比寫兩個方向各算一遍更清楚。

\begin{code}{c}{若 a 大於 b 就先交換}
if (a > b) {
    int t = a;
    a = b;
    b = t;
}
for (int i = a; i <= b; i++) {
    sum += i;
}
\end{code}

範例 \thechapter.39 除了印總和，還把算式 $a + (a+1) + \cdots = s$ 完整
輸出。 迴圈裡用旗標 \texttt{first} 控制加號，避免算式開頭多一個
\texttt{+}。

\begin{example}{c}{(sum-equation.c) 連續整數求和並印出算式}
#include <stdio.h>

int main(void) {
    int a, b;
    scanf("%d %d", &a, &b);

    if (a > b) {
        int t = a;
        a = b;
        b = t;
    }

    long long sum = 0;
    int first = 1;
    for (int i = a; i <= b; i++) {
        if (!first) {
            printf("+");
        }
        printf("%d", i);
        sum += i;
        first = 0;
    }
    printf("=%lld\n", sum);
    return 0;
}
\end{example}

\iosample{
1 5 \\
}{
1+2+3+4+5=15 \\
}

\keypoint{連續整數求和雖然有公式 $(A+B)(B-A+1)/2$，但用迴圈實作更直覺，
也為之後「不是等差數列、不能套公式」的累加題打好基礎。}

\begin{drills}[label=drills:loop-range-sum]
\item 如果輸入的 $a$ 大於 $b$，程式在開始累加前會先做什麼？
\item \texttt{sum-equation.c} 裡的旗標 \texttt{first} 是用來避免什麼
問題？
\end{drills}

\section{累加的變奏：次方和、下一個倍數、最小公倍數}
\label{sec:loop-accum-variants}

同一層 \texttt{for} 可以同時累加好幾種量。 範例 \thechapter.40 讀
\texttt{n}，一次算出 $1$ 到 $n$ 的和、平方和、立方和。 $n = 4$ 時是
$10$、$30$、$100$。

\begin{example}{c}{(sum-powers.c) 連續整數的和、平方和、立方和}
#include <stdio.h>

int main(void) {
    int n;
    scanf("%d", &n);
    int sum = 0;
    int sq = 0;
    int cu = 0;
    for (int i = 1; i <= n; i++) {
        sum += i;
        sq += i * i;
        cu += i * i * i;
    }
    printf("%d %d %d\n", sum, sq, cu);
    return 0;
}
\end{example}

\iosample{
4 \\
}{
10 30 100 \\
}

有些 \texttt{while} 不是在「數次數」，而是在「把狀態推到下一個合格
點」。 範例 \thechapter.41 從 $0$ 每次加 $5$，直到不小於輸入的
\texttt{j}。 也就是「$5$ 的倍數裡，不小於 $j$ 的最小那個」。
\texttt{j = 17} 得到 $20$；\texttt{j = 15} 得到 $15$。

\begin{example}{c}{(next-five.c) 不小於 j 的最小 5 的倍數}
#include <stdio.h>

int main(void) {
    int j;
    scanf("%d", &j);
    int i = 0;
    while (i < j) {
        i += 5;
    }
    printf("%d\n", i);
    return 0;
}
\end{example}

\iosample{
17 \\
}{
20 \\
}

最小公倍數可以從已經會的輾轉相除再走一步：$\mathrm{lcm}(a,b) =
ab / \gcd(a,b)$。 為了少溢位，先除後乘：\texttt{a / gcd * b}。 範例
\thechapter.42 對 \texttt{12 18} 印出 \texttt{36}。

\begin{example}{c}{(lcm-two.c) 用 gcd 迴圈求最小公倍數}
#include <stdio.h>

int main(void) {
    int a, b;
    scanf("%d %d", &a, &b);
    int x = a;
    int y = b;
    while (y != 0) {
        int t = x % y;
        x = y;
        y = t;
    }
    printf("%d\n", a / x * b);
    return 0;
}
\end{example}

\iosample{
12 18 \\
}{
36 \\
}

\begin{drills}[label=drills:loop-accum-variants]
\item \texttt{next-five.c} 輸入 \texttt{j = 17}，\texttt{i} 從 $0$
開始每次加 $5$，最後印出多少？
\item \texttt{lcm-two.c} 怎麼從 \texttt{gcd} 算出最小公倍數？
\item 為什麼要寫成 \texttt{a / gcd * b} 而不是 \texttt{a * b / gcd}？
\end{drills}

\section{位數處理：剝除與反轉}
\label{sec:loop-digits}

整數在記憶體裡是以二進位存的，但我們常需要按\textbf{十進位}的個位、十位
來處理。 取餘數 \texttt{n \% 10} 得到最後一位，再 \texttt{n /= 10} 把
最後一位「剝掉」。 這個 \texttt{while (n > 0)} 模式在數論題目裡非常常見。

\begin{example}{c}{(digit-sum.c) 各位數字相加}
#include <stdio.h>

int main(void) {
    int n;
    scanf("%d", &n);

    int s = 0;
    while (n > 0) {
        s += n % 10;
        n /= 10;
    }
    printf("%d\n", s);
    return 0;
}
\end{example}

\iosample{
12345 \\
}{
15 \\
}

範例 \thechapter.44 用同一套剝除邏輯，但把每一位「接」到新整數的末尾，
達成數字反轉。

\begin{example}{c}{(reverse-digits.c) 反轉整數的十進位數字}
#include <stdio.h>

int main(void) {
    int n;
    scanf("%d", &n);

    int rev = 0;
    while (n > 0) {
        rev = rev * 10 + n % 10;
        n /= 10;
    }
    printf("%d\n", rev);
    return 0;
}
\end{example}

\iosample{
12345 \\
}{
54321 \\
}

以 $12345$ 走一遍反轉：\texttt{rev} 依序變成 $5 \to 54 \to 543 \to
5432 \to 54321$。 每次 \texttt{rev * 10} 把舊結果左移一位，再加上新的
個位。

\stylepoint{位數處理的 \texttt{while (n > 0)} 迴圈，每一輪只做兩件事：
\texttt{n \% 10} 取個位、\texttt{n /= 10} 剝掉個位。 把這個節奏記熟，
很多題目都能套同一個骨架。}

\begin{drills}[label=drills:loop-digits]
\item \texttt{digit-sum.c} 用哪兩個運算子，把整數的個位剝出來、再把
它去掉？
\item \texttt{reverse-digits.c} 裡 \texttt{rev = rev * 10 + n \% 10;}
這一行在做什麼？
\end{drills}

\section{累乘、階乘型乘積與狀態更新}
\label{sec:loop-product}

連加對應的是 \texttt{sum += ...}，連乘則是 \texttt{product *= ...}。
排列數 $P(N,R) = N \times (N-1) \times \cdots \times (N-R+1)$ 需要
\texttt{long long}，因為乘積很快超過 \texttt{int} 範圍。

\begin{example}{c}{(falling-factorial.c) 計算 P(N,R)}
#include <stdio.h>

int main(void) {
    int n, r;
    scanf("%d %d", &n, &r);

    long long ans = 1;
    for (int i = 0; i < r; i++) {
        ans *= (n - i);
    }
    printf("%lld\n", ans);
    return 0;
}
\end{example}

\iosample{
5 3 \\
}{
60 \\
}

有些 \texttt{while} 迴圈每一輪要同時更新好幾個變數。 範例 \thechapter.46
模擬細菌：每天族群乘以 $3$，直到達到目標。

\begin{example}{c}{(bacteria-growth.c) 三倍成長直到達標}
#include <stdio.h>

int main(void) {
    int pop, target;
    scanf("%d %d", &pop, &target);

    int days = 0;
    while (pop < target) {
        pop *= 3;
        days++;
    }
    printf("%d\n", days);
    return 0;
}
\end{example}

\iosample{
1 27 \\
}{
3 \\
}

範例 \thechapter.47 則是「兌換迴圈」：每 $5$ 根棍子可換 $1$ 根，換來的
又能繼續湊。 迴圈裡同時更新「目前手上」與「累計總數」兩個狀態。

\begin{example}{c}{(exchange-sticks.c) 每 5 根兌換 1 根}
#include <stdio.h>

int main(void) {
    int sticks;
    scanf("%d", &sticks);

    int total = sticks;
    while (sticks >= 5) {
        int gained = sticks / 5;
        total += gained;
        sticks = sticks % 5 + gained;
    }
    printf("%d\n", total);
    return 0;
}
\end{example}

\iosample{
23 \\
}{
28 \\
}

$23$ 根：先換 $4$ 根（總 $27$，剩 $7$），再換 $1$ 根（總 $28$，剩 $3$），
$3 < 5$ 結束。

\keypoint{累乘、累加一樣要選夠大的型別；乘積比加法更容易溢出，看到排列、
組合、階乘這類關鍵字，優先考慮 \texttt{long long}。}

\begin{drills}[label=drills:loop-product]
\item \texttt{falling-factorial.c} 為什麼用 \texttt{long long} 存
\texttt{ans}，而不是 \texttt{int}？
\item \texttt{bacteria-growth.c} 的 \texttt{while} 迴圈裡，每一輪同時
更新哪兩個變數？
\end{drills}

\section{歐幾里得算法與質數判定}
\label{sec:loop-gcd-prime}

\textbf{輾轉相除法}（Euclid 算法）用「大數除以小數、取餘數、把小數當新
的大數」反覆進行，直到餘數為 $0$，此時的小數就是最大公因數。 整個過程
就是一個 \texttt{while (b != 0)} 迴圈。

\begin{example}{c}{(gcd-euclid.c) 輾轉相除法求最大公因數}
#include <stdio.h>

int main(void) {
    int a, b;
    scanf("%d %d", &a, &b);

    while (b != 0) {
        int t = a % b;
        a = b;
        b = t;
    }
    printf("%d\n", a);
    return 0;
}
\end{example}

\iosample{
48 18 \\
}{
6 \\
}

判定質數則用 \texttt{for} 從 $2$ 試除到 $\sqrt{n}$。 只要找到任何一個
因數就可以用前面學過的 \texttt{break} 提早結束，不必試完所有數字。

\begin{example}{c}{(prime-check.c) 試除法判定質數}
#include <stdio.h>

int main(void) {
    int n;
    scanf("%d", &n);

    if (n < 2) {
        printf("NO\n");
        return 0;
    }

    int isPrime = 1;
    for (int d = 2; (long long)d * d <= n; d++) {
        if (n % d == 0) {
            isPrime = 0;
            break;
        }
    }
    printf("%s\n", isPrime ? "YES" : "NO");
    return 0;
}
\end{example}

\iosample{
17 \\
}{
YES \\
}

\keypoint{試除到 $\sqrt{n}$ 就夠，因為若 $n = ab$ 且 $a \le b$，則
$a \le \sqrt{n}$；用 \texttt{(long long)d * d <= n} 避免 \texttt{d * d}
在 \texttt{d} 很大時溢出。}

\begin{drills}[label=drills:loop-gcd-prime]
\item 輾轉相除法裡，\texttt{while (b != 0)} 迴圈結束時，哪個變數存的
就是最大公因數？
\item \texttt{prime-check.c} 為什麼只試除到 $\sqrt{n}$ 就夠，不用試到
$n$？
\item 條件寫成 \texttt{(long long)d * d <= n} 而不是 \texttt{d * d <= n}
的目的是什麼？
\end{drills}

\section{Collatz 序列}
\label{sec:loop-collatz}

$3n+1$ 問題（Collatz 猜想）：若 $n$ 是偶數就除以 $2$，奇數就變成
$3n+1$，重複直到變成 $1$。 沒人證明是否一定會到 $1$，但程式可以忠實
模擬這個過程。

範例 \thechapter.50 印出整段序列；範例 \thechapter.51 只數長度（包含起
點）。 兩者都適合用 \texttt{while (n != 1)}，差別在要不要把中間每一步
都印出來。

\begin{example}{c}{(collatz-sequence.c) 印出 Collatz 序列}
#include <stdio.h>

int main(void) {
    int n;
    scanf("%d", &n);

    int first = 1;
    while (1) {
        if (!first) {
            printf(" ");
        }
        printf("%d", n);
        first = 0;
        if (n == 1) {
            break;
        }
        if (n % 2 == 0) {
            n /= 2;
        } else {
            n = 3 * n + 1;
        }
    }
    printf("\n");
    return 0;
}
\end{example}

\iosample{
6 \\
}{
6 3 10 5 16 8 4 2 1 \\
}

\begin{example}{c}{(collatz-length.c) 計算 Collatz 序列長度}
#include <stdio.h>

int main(void) {
    int n;
    scanf("%d", &n);

    int len = 1;
    while (n != 1) {
        if (n % 2 == 0) {
            n /= 2;
        } else {
            n = 3 * n + 1;
        }
        len++;
    }
    printf("%d\n", len);
    return 0;
}
\end{example}

\iosample{
6 \\
}{
9 \\
}

\keypoint{Collatz 迴圈裡，奇數步 \texttt{3 * n + 1} 可能讓 \texttt{n}
短時間內變很大；若題目上限很大，要確認 \texttt{int} 是否夠用，必要時
改用 \texttt{long long}。}

\begin{drills}[label=drills:loop-collatz]
\item Collatz 規則：$n$ 是偶數時怎麼變？是奇數時怎麼變？
\item \texttt{collatz-length.c} 跟 \texttt{collatz-sequence.c} 用的是
同一套規則，兩支程式的輸出差在哪裡？
\end{drills}

\section{因數和、餘數分類與同號連續段}
\label{sec:loop-divisor-classify}

把 $1$ 到 $n-1$ 中所有整除 $n$ 的因數加起來，可以判斷 $n$ 是完數
（Perfect）、盈數（Abundant）還是虧數（Deficient）。 範例 \thechapter.52
用一個 \texttt{for} 累加真因數和。

\begin{example}{c}{(classify-divisors.c) 依真因數和分類整數}
#include <stdio.h>

int main(void) {
    int n;
    scanf("%d", &n);

    int sum = 0;
    for (int d = 1; d < n; d++) {
        if (n % d == 0) {
            sum += d;
        }
    }

    if (sum == n) {
        printf("Perfect\n");
    } else if (sum > n) {
        printf("Abundant\n");
    } else {
        printf("Deficient\n");
    }
    return 0;
}
\end{example}

\iosample{
28 \\
}{
Perfect \\
}

範例 \thechapter.53 讀 $n$ 個整數，依除以 $3$ 的餘數分三類計數。 這
跟陣列章的「分桶」概念類似，但這裡只用三個計數器，不必開陣列。

\begin{example}{c}{(residue-count.c) 依 mod 3 餘數分類計數}
#include <stdio.h>

int main(void) {
    int n;
    scanf("%d", &n);

    int c0 = 0, c1 = 0, c2 = 0;
    for (int i = 0; i < n; i++) {
        int x;
        scanf("%d", &x);
        if (x % 3 == 0) {
            c0++;
        } else if (x % 3 == 1) {
            c1++;
        } else {
            c2++;
        }
    }
    printf("%d %d %d\n", c0, c1, c2);
    return 0;
}
\end{example}

\iosample{
5 \\
3 4 5 6 7 \\
}{
2 2 1 \\
}

範例 \thechapter.54 處理「連續同號整數的長度」。 掃描時若符號跟上一個
相同就延長計數，否則輸出上一段並重新開始。 這是\textbf{run-length}（游
程編碼）的雛形，之後字串、序列題目還會再遇到。

\begin{example}{c}{(run-length-sign.c) 同號連續段長度}
#include <stdio.h>

int main(void) {
    int n;
    scanf("%d", &n);

    int prev, count = 0, first = 1;
    for (int i = 0; i < n; i++) {
        int x;
        scanf("%d", &x);
        if (i == 0) {
            prev = x;
            count = 1;
            continue;
        }
        if ((prev > 0 && x > 0) || (prev < 0 && x < 0)) {
            count++;
        } else {
            if (!first) {
                printf(" ");
            }
            printf("%d", (prev > 0 ? 1 : -1) * count);
            first = 0;
            prev = x;
            count = 1;
        }
    }

    if (n > 0) {
        if (!first) {
            printf(" ");
        }
        printf("%d\n", (prev > 0 ? 1 : -1) * count);
    }
    return 0;
}
\end{example}

\iosample{
6 \\
1 2 3 -1 -2 5 \\
}{
3 -2 1 \\
}

正數段 $1,2,3$ 長度 $3$ 輸出 \texttt{3}；負數段 $-1,-2$ 長度 $2$ 輸出
\texttt{-2}（符號代表正負）；最後單獨的 $5$ 輸出 \texttt{1}。

\begin{drills}[label=drills:loop-divisor-classify]
\item \texttt{classify-divisors.c} 怎麼判斷一個數是不是「完數」
（Perfect）？
\item \texttt{residue-count.c} 用了幾個計數器？分別對應除以 $3$ 餘幾
的情況？
\item \texttt{run-length-sign.c} 在處理的是哪一種常見模式（跟後面字
串章節會再遇到的哪個概念一樣）？
\end{drills}

\section{日期計算：從年月日推算星期}
\label{sec:loop-day-of-week}

日期題目通常要處理：各月天數不同、閏年二月多一天、從某個基準日累加總天
數再取餘數。 範例 \thechapter.55 簡化版：假設該年 1 月 1 日是星期五
（編號 $6$，若星期日為 $0$），讀入月日後累加從年初到該日前的天數。

\begin{example}{c}{(day-of-week.c) 由月日推算星期（簡化基準）}
#include <stdio.h>

int isLeap(int y) {
    return (y % 400 == 0) || (y % 4 == 0 && y % 100 != 0);
}

int main(void) {
    int y, m, d;
    scanf("%d %d %d", &y, &m, &d);

    int days[] = {0, 31, 28, 31, 30, 31, 30, 31, 31, 30, 31, 30, 31};
    if (isLeap(y)) {
        days[2] = 29;
    }

    int total = 0;
    for (int i = 1; i < m; i++) {
        total += days[i];
    }
    total += d - 1;

    int w = (total + 6) % 7;
    printf("%d\n", w);
    return 0;
}
\end{example}

\iosample{
2024 3 1 \\
}{
3 \\
}

閏年判斷 \texttt{isLeap} 抽出成函式，主程式用 \texttt{for} 累加整月天
數。 這種「查表 + 累加 + 取餘」的組合，是日期類題目的共通骨架。 更完整
的月曆輸出（格子排版、 weekday offset）只是把同一套天數計算嵌進巢狀
迴圈裡控制列與欄。

\keypoint{日期題目裡，閏年規則（四年一閏、百年不閏、四百年再閏）跟「各
月天數表」一定要先寫對；邏輯錯在這裡，後面再漂亮的迴圈也救不了。}

\begin{drills}[label=drills:loop-day-of-week]
\item \texttt{isLeap} 函式判斷閏年的三個條件是什麼？
\item 為什麼程式要在確定是不是閏年之後，才把 \texttt{days[2]} 改成
$29$？
\end{drills}

\section{進階模擬：多狀態與複合迴圈}
\label{sec:loop-advanced-sim}

題庫後段（分數 $80$ 以上）的題目，往往把多種技巧疊在一起：在網格上模擬
機器人移動、光線在鏡面間反彈、多個容器之間反覆倒水直到平衡、或是在狀態
空間裡用 \texttt{while} 迭代直到收斂。 這些題目沒有單一「萬用片語」，
但共通結構是：

\begin{enumerate}
    \item 用變數（或陣列）完整描述「系統現狀」；
    \item 外層 \texttt{while} 控制「模擬還要不要繼續」；
    \item 內層 \texttt{for} 處理「這一輪裡每個物件各做一步」；
    \item 每步更新後檢查終止條件（碰撞、填滿、空池、超過步數上限等）。
\end{enumerate}

練習時建議先在紙上對小測資手工模擬幾輪，確認「每一輪哪些變數會變、什麼
情況該停」，再把它翻成程式。 若單純用 \texttt{for} 掃一次無法表達「要
重複調整直到穩定」，就該懷疑需要外層 \texttt{while}。 這是從基礎迴圈
過渡到模擬題的分水嶺。

題庫裡像「兩球在格子裡彈跳並交換位置」「多個杯子與無限水池之間倒水」
「機器人找燃料並在環狀地圖上相遇」這類高分題，都建立在本章前面所有
技巧之上：輸入驗證、哨兵、多筆測資、巢狀迴圈、\texttt{break}、以及多
變數狀態更新。 把這一章的基礎練熟，再回頭挑戰那些題目，會比一開始就
硬啃模擬題有效率得多。

\begin{drills}[label=drills:loop-advanced-sim]
\item 進階模擬題的共通結構列出的四個步驟依序是什麼？
\item 當「單純用 \texttt{for} 掃一次無法表達」時，通常代表需要哪一種
額外的迴圈結構？
\end{drills}

\addcontentsline{toc}{section}{練習解答}
\section*{練習解答}

各節練習的答案收在這裡。 短題只給能一句話核對的結果。

\paragraph{第~\ref{sec:loop-for-basics}~節}
\begin{enumerate}
\item 一次都不會執行：條件 \texttt{0 < 0} 一開始就不成立。
\item 不能，\texttt{i} 的作用域只到這一句 \texttt{for} 結束為止。
\item \texttt{i < n} 跑 $n$ 次；\texttt{i <= n} 跑 $n + 1$ 次。
\end{enumerate}

\paragraph{第~\ref{sec:loop-countdown}~節}
\begin{enumerate}
\item \texttt{i >= 0} 那個會多印一個；多印的是 $0$。
\item 因為最後一個數字（$0$）後面不該有空白；\texttt{if (i > 0)} 讓
只有非最後一個數字才補空白。
\item 每個數字後面都會多印一個空白，行尾會多出一個空格。
\end{enumerate}

\paragraph{第~\ref{sec:loop-while-basics}~節}
\begin{enumerate}
\item 跑三輪：$24 \to 12 \to 6 \to 3$，$3$ 是奇數時結束。
\item 條件永遠成立（\texttt{n} 不會變），形成無窮迴圈。
\item 開始。
\end{enumerate}

\paragraph{第~\ref{sec:loop-for-while-equiv}~節}
\begin{enumerate}
\item 放在本體的最後面（先執行本體，再更新，才回頭測條件）。
\item 都是 $0$ 次。
\item $0$、$1$、$2$。
\end{enumerate}

\paragraph{第~\ref{sec:loop-do-while}~節}
\begin{enumerate}
\item \texttt{do-while} 保證本體至少執行一次；\texttt{while} 若一開始
條件就不成立，可能一次都不執行。
\item 印出 \texttt{5}：因為本體在條件被檢查之前就先跑過一次。
\item 一次：讀到正數後條件 \texttt{x <= 0} 不成立，迴圈立刻結束。
\end{enumerate}

\paragraph{第~\ref{sec:loop-break-continue}~節}
\begin{enumerate}
\item \texttt{break} 整段離開迴圈；\texttt{continue} 只結束這一輪，
接著跑下一輪。
\item 印出 \texttt{-1}：$1$ 到 $5$ 裡沒有 $7$ 的倍數。
\item 跳過的是 \texttt{-2}；\texttt{sum = 7 + 3 + 8 = 18}。
\end{enumerate}

\paragraph{第~\ref{sec:loop-two-stage}~節}
\begin{enumerate}
\item 避免 $n$ 被印兩次——第一段已經印過 $n$ 了。
\item $n$ 會被連續印兩次（一次在第一段結尾，一次在第二段開頭）。
\end{enumerate}

\paragraph{第~\ref{sec:loop-nested-shapes}~節}
\begin{enumerate}
\item 外層通常控制「第幾列」，內層控制「這一列印幾次或印什麼」。
\item 印 $i$ 次；第 $3$ 列印出 \texttt{333}。
\item \texttt{number-triangle.c} 每列重複同一個數字；
\texttt{repeated-pattern.c} 每列印遞增序列 $1$ 到 $i$。
\end{enumerate}

\paragraph{第~\ref{sec:loop-mult-table}~節}
\begin{enumerate}
\item 把每個數字用固定寬度（$2$ 格）對齊輸出，讓欄位排整齊。
\item $j = 0, 3$ 時 $(0+j) \% 3 == 0$ 成立，也就是第 $0$、$3$ 欄印
\texttt{@}。
\end{enumerate}

\paragraph{第~\ref{sec:loop-pyramid}~節}
\begin{enumerate}
\item 先印 $n - i$ 個底線，再印 $i$ 個數字。
\item 第一段負責左邊的留白（底線），第二段負責這一列的內容（數字）。
\end{enumerate}

\paragraph{第~\ref{sec:loop-input-validation}~節}
\begin{enumerate}
\item 因為一定要先讀一次才能判斷合不合法，\texttt{do-while} 保證本體
至少執行一次。
\item 跳過輸入裡殘留的空白字元（包含前一次讀取留下的換行）。
\item 三次；最後印出 \texttt{85}。
\end{enumerate}

\paragraph{第~\ref{sec:loop-sentinel}~節}
\begin{enumerate}
\item 一個約定好、代表「資料結束了」的特殊輸入值。
\item \texttt{-1} 那版不把哨兵本身計入結果；\texttt{0} 那版連最後讀到
的 $0$ 都一起加進總和。
\item 依序印出 \texttt{AGAIN}、\texttt{AGAIN}、\texttt{BYE}。
\end{enumerate}

\paragraph{第~\ref{sec:loop-multi-case}~節}
\begin{enumerate}
\item 外層用 \texttt{for (int i = 0; i < n; i++)} 包住處理每一組資料
的邏輯。
\item \texttt{case-prefix.c} 多印了 \texttt{Case k:} 這個前綴。
\item \texttt{while (scanf(...) != EOF)}。
\end{enumerate}

\paragraph{第~\ref{sec:loop-range-sum}~節}
\begin{enumerate}
\item 先把 $a$、$b$ 交換，確保 $a \le b$ 再累加。
\item 避免算式開頭多印一個 \texttt{+} 號。
\end{enumerate}

\paragraph{第~\ref{sec:loop-accum-variants}~節}
\begin{enumerate}
\item $20$（$5$ 的倍數裡不小於 $17$ 的最小值）。
\item 先用 \texttt{while} 迴圈算出 \texttt{gcd}，再用
\texttt{a / gcd * b}。
\item 避免 \texttt{a * b} 在中間步驟就超過型別能表示的範圍（先除可以
讓中間值小一點）。
\end{enumerate}

\paragraph{第~\ref{sec:loop-digits}~節}
\begin{enumerate}
\item \texttt{\% 10} 取出個位，\texttt{/= 10} 把個位剝掉。
\item 把新讀到的個位「接」到 \texttt{rev} 的末尾，達成反轉的效果。
\end{enumerate}

\paragraph{第~\ref{sec:loop-product}~節}
\begin{enumerate}
\item 因為排列數的乘積很快就會超過 \texttt{int} 能表示的範圍。
\item 同時更新族群數量 \texttt{pop}（乘以 $3$）與天數 \texttt{days}
（加 $1$）。
\end{enumerate}

\paragraph{第~\ref{sec:loop-gcd-prime}~節}
\begin{enumerate}
\item \texttt{a}：迴圈結束時 \texttt{a} 就是最大公因數。
\item 因為若 $n = a \times b$ 且 $a \le b$，則必有 $a \le \sqrt{n}$，
超過 $\sqrt{n}$ 不會再找到新的因數。
\item 避免 \texttt{d * d} 在 \texttt{d} 很大時，用 \texttt{int} 運算會
先溢位。
\end{enumerate}

\paragraph{第~\ref{sec:loop-collatz}~節}
\begin{enumerate}
\item 偶數除以 $2$；奇數變成 $3n + 1$。
\item \texttt{collatz-length.c} 只印最終的步數；\texttt{collatz-sequence.c}
把過程中每一步的值都印出來。
\end{enumerate}

\paragraph{第~\ref{sec:loop-divisor-classify}~節}
\begin{enumerate}
\item 真因數和 \texttt{sum} 剛好等於 $n$ 本身時就是完數。
\item 三個計數器 \texttt{c0}、\texttt{c1}、\texttt{c2}，分別對應除以
$3$ 餘 $0$、$1$、$2$ 的情況。
\item 游程編碼（run-length encoding）：把連續相同（這裡是同號）的一
段資料記成「值＋長度」。
\end{enumerate}

\paragraph{第~\ref{sec:loop-day-of-week}~節}
\begin{enumerate}
\item 能被 $400$ 整除；或者能被 $4$ 整除但不能被 $100$ 整除。
\item 因為閏年二月是 $29$ 天、平年是 $28$ 天，要先判斷完閏年才知道
\texttt{days[2]} 該填哪個值。
\end{enumerate}

\paragraph{第~\ref{sec:loop-advanced-sim}~節}
\begin{enumerate}
\item 描述系統現狀、外層 \texttt{while} 控制模擬要不要繼續、內層
\texttt{for} 處理每個物件各走一步、每步後檢查終止條件。
\item 通常代表需要外層 \texttt{while}（因為要「重複調整直到穩定」，
次數不是事先就知道的）。
\end{enumerate}
